\documentclass[11pt,a4paper]{article}

\usepackage[utf8]{inputenc}
\usepackage[T1]{fontenc}
\usepackage[english]{babel}
\usepackage{amsmath,amssymb,mathtools}
\usepackage{geometry}

\geometry{margin=2.6cm}

\newcommand{\R}{\mathbb{R}}
\newcommand{\norm}[1]{\left\lVert #1\right\rVert}

\title{Computation of Multipliers Using Broyden's Method\\
for a Generic Numerical Integrator}
\author{}
\date{}

\begin{document}

\maketitle

\section{Objective}

In many numerical methods, the advance from $z_n$ to $z_{n+1}$ depends on a
multiplier or correction vector
\[
    \mu\in\R^q
\]
that must be chosen so that the result satisfies a coupling, projection, or
constraint condition. The computation of $\mu$ leads to a nonlinear system of
the form
\begin{equation}
    R_n(\mu)=0,
    \qquad
    R_n:\R^q\longrightarrow\R^q.
    \label{eq:generic-root-problem}
\end{equation}

Broyden's method makes it possible to solve \eqref{eq:generic-root-problem}
without recomputing the exact Jacobian of the residual at every iteration.
Its formulation depends only on the ability to evaluate $R_n(\mu)$ and not on
the internal structure of the integrator. Therefore, it can be used with
ABBA, BM4, or any other numerical method.

\section{Generic Construction of the Residual}

We separate the procedure into four maps.

\subsection{Input map or lifting}

The physical state $z_n\in\R^d$ and the multiplier $\mu\in\R^q$ are
transformed into the state on which the numerical method acts:
\begin{equation}
    Y_n(\mu)=\mathcal E_n(z_n,\mu)\in\R^m.
    \label{eq:lifting}
\end{equation}
The map $\mathcal E_n$ may duplicate the state, shift some of its components,
or introduce auxiliary variables.

\subsection{Base integrator}

Let
\begin{equation}
    \Psi_{h,n}:\R^m\longrightarrow\R^m
    \label{eq:base-integrator}
\end{equation}
be the one-step map of the chosen integrator. The state obtained after
applying the method is
\begin{equation}
    Y_{n+1}(\mu)
    =\Psi_{h,n}\bigl(Y_n(\mu)\bigr)
    =\Psi_{h,n}\bigl(\mathcal E_n(z_n,\mu)\bigr).
    \label{eq:numerical-output}
\end{equation}
The notation $\Psi_{h,n}$ represents, in an abstract way, any scheme: ABBA,
BM4, another composition method, or even an integrator not based on
compositions.

\subsection{Compatibility condition}

The result must satisfy a condition written through
\begin{equation}
    \mathcal C_n:\R^m\times\R^q\longrightarrow\R^q.
    \label{eq:compatibility-map}
\end{equation}
The associated residual is defined by
\begin{equation}
    \boxed{
    R_n(\mu)
    :=\mathcal C_n
    \left(
        \Psi_{h,n}\bigl(\mathcal E_n(z_n,\mu)\bigr),
        \mu
    \right).}
    \label{eq:generic-residual}
\end{equation}
Thus, finding the multiplier is equivalent to solving
\[
    R_n(\mu_n)=0.
\]

The function $\mathcal C_n$ may represent, among other possibilities, the
equality between several copies of the same variable, a constraint condition
$g(z)=0$, or a projection equation.

\subsection{Reconstruction of the physical state}

Once $\mu_n$ has been found, the physical state is reconstructed through a
map
\begin{equation}
    z_{n+1}
    =\mathcal P_n\bigl(Y_{n+1}(\mu_n),\mu_n\bigr),
    \label{eq:reconstruction}
\end{equation}
where $\mathcal P_n$ depends on the specific architecture of the method.

In this way, the part specific to each integrator is contained in
\[
    \mathcal E_n,
    \qquad
    \Psi_{h,n},
    \qquad
    \mathcal C_n,
    \qquad
    \mathcal P_n,
\]
whereas Broyden's algorithm remains unchanged.

\section{Exact Jacobian of the Residual}

Newton's method applied to \eqref{eq:generic-root-problem} would be
\begin{equation}
    \mu^{(k+1)}
    =\mu^{(k)}
    -K\bigl(\mu^{(k)}\bigr)^{-1}
     R_n\bigl(\mu^{(k)}\bigr),
    \label{eq:newton}
\end{equation}
where
\begin{equation}
    K(\mu):=D_\mu R_n(\mu)\in\R^{q\times q}.
    \label{eq:residual-jacobian}
\end{equation}

If all maps are differentiable, the chain rule applied to
\eqref{eq:generic-residual} gives
\begin{equation}
\boxed{
\begin{aligned}
    K(\mu)
    ={}&D_Y\mathcal C_n\bigl(Y_{n+1}(\mu),\mu\bigr)
       \,D\Psi_{h,n}\bigl(Y_n(\mu)\bigr)
       \,D_\mu\mathcal E_n(z_n,\mu)\\
      &+D_\mu\mathcal C_n\bigl(Y_{n+1}(\mu),\mu\bigr).
\end{aligned}}
\label{eq:generic-jacobian}
\end{equation}

This is the Jacobian of the nonlinear problem used to compute $\mu$. It
should not be confused with the physical Jacobian of the step
$z_n\mapsto z_{n+1}$.

Computing \eqref{eq:generic-jacobian} exactly at every iteration may be
expensive because it contains the derivative of the entire integrator,
$D\Psi_{h,n}$. Broyden's method replaces this matrix by an approximation that
is updated using successive evaluations of the residual.

\section{Broyden's Method}

Suppose that the current iterate $\mu^{(k)}$, the residual
\begin{equation}
    R_k:=R_n\bigl(\mu^{(k)}\bigr)
    \label{eq:current-residual}
\end{equation}
and an approximation
\[
    K_k\simeq D_\mu R_n\bigl(\mu^{(k)}\bigr)
\]
are known.

The correction $\Delta\mu_k$ is obtained by solving
\begin{equation}
    K_k\Delta\mu_k=-R_k,
    \label{eq:linear-correction}
\end{equation}
and the multiplier is updated as
\begin{equation}
    \mu^{(k+1)}=\mu^{(k)}+\Delta\mu_k.
    \label{eq:multiplier-update}
\end{equation}

The complete integrator is then evaluated again in order to compute
\[
    R_{k+1}:=R_n\bigl(\mu^{(k+1)}\bigr).
\]
Define
\begin{equation}
    s_k:=\mu^{(k+1)}-\mu^{(k)},
    \qquad
    y_k:=R_{k+1}-R_k.
    \label{eq:secant-vectors}
\end{equation}

The vector $s_k$ is the change introduced in the multiplier, and $y_k$ is the
observed change in the residual. The new Jacobian approximation must
incorporate this information through the secant condition
\begin{equation}
    K_{k+1}s_k=y_k.
    \label{eq:secant-condition}
\end{equation}

The good Broyden update is given by
\begin{equation}
    \boxed{
    K_{k+1}
    =K_k+
    \frac{\bigl(y_k-K_ks_k\bigr)s_k^{\mathsf T}}
         {s_k^{\mathsf T}s_k}.}
    \label{eq:good-broyden}
\end{equation}
Indeed,
\[
\begin{aligned}
    K_{k+1}s_k
    &=K_ks_k
      +\frac{\bigl(y_k-K_ks_k\bigr)s_k^{\mathsf T}s_k}
             {s_k^{\mathsf T}s_k}\\
    &=y_k.
\end{aligned}
\]
Therefore, \eqref{eq:good-broyden} is a multidimensional generalization of
the secant method. Moreover, it is the smallest modification of $K_k$ in the
Frobenius norm that satisfies \eqref{eq:secant-condition}.

\section{Initialization}

The choice of the initial values depends on the specific problem, not on
Broyden's algorithm itself.

\subsection{Initial value of the multiplier}

A common choice is
\begin{equation}
    \mu^{(0)}=0,
    \label{eq:zero-initial-multiplier}
\end{equation}
when $\mu$ represents a small correction and the uncorrected method provides a
sufficiently close approximation. This choice does not state that $\mu=0$ is
the solution; it only fixes the first point at which $R_n$ is evaluated.

The multiplier from the previous time step may also be used,
\[
    \mu^{(0)}=\mu_{n-1},
\]
or any available predictor. In every case, Broyden's method starts by
evaluating the corresponding residual exactly.

\subsection{Initial matrix}

The matrix $K_0$ should approximate
\[
    D_\mu R_n\bigl(\mu^{(0)}\bigr).
\]
It can be obtained in different ways:

\begin{enumerate}
    \item by evaluating the exact Jacobian only at the beginning;

    \item by using an asymptotic approximation valid for small $h$;

    \item by finite differences, computing its columns as
    \begin{equation}
        K_0e_j
        \approx
        \frac{R_n(\mu^{(0)}+\varepsilon e_j)
              -R_n(\mu^{(0)})}{\varepsilon};
        \label{eq:finite-difference-initialization}
    \end{equation}

    \item by using a simple constant matrix, for example
    $K_0=\alpha I_q$, when the structure of the problem justifies this
    approximation.
\end{enumerate}

The choice $K_0=4I_q$ is not part of Broyden's method in general. It is an
initialization specific to certain formulations with two symmetric copies.
For a different lifting, a different compatibility condition, or even a
different residual, the corresponding approximation must be used.

\section{Formulation Using the Approximate Inverse}

Instead of storing $K_k$, one may directly store
\[
    H_k\simeq K_k^{-1}.
\]
The iteration can then be written as
\begin{equation}
    \mu^{(k+1)}=\mu^{(k)}-H_kR_k.
    \label{eq:inverse-step}
\end{equation}
The inverse update associated with \eqref{eq:good-broyden} is
\begin{equation}
    \boxed{
    H_{k+1}
    =H_k+
    \frac{\bigl(s_k-H_ky_k\bigr)s_k^{\mathsf T}H_k}
         {s_k^{\mathsf T}H_ky_k},}
    \label{eq:inverse-broyden}
\end{equation}
provided that the denominator is nonzero. The initial matrix must satisfy
\[
    H_0\simeq K_0^{-1}.
\]

\section{Generic Algorithm for One Time Step}

To advance from $z_n$ to $z_{n+1}$, the following procedure is applied:

\begin{enumerate}
    \item Choose $\mu^{(0)}$ and an initial approximation $K_0$.

    \item For $k=0,1,2,\ldots$, construct
    \[
        Y_n^{(k)}=\mathcal E_n\bigl(z_n,\mu^{(k)}\bigr).
    \]

    \item Perform one complete step of the chosen integrator:
    \[
        Y_{n+1}^{(k)}=\Psi_{h,n}\bigl(Y_n^{(k)}\bigr).
    \]

    \item Evaluate the residual
    \[
        R_k
        =\mathcal C_n\bigl(Y_{n+1}^{(k)},\mu^{(k)}\bigr).
    \]

    \item If
    \[
        \norm{R_k}_2\leq\mathrm{tol},
    \]
    stop the iterations and set $\mu_n=\mu^{(k)}$.

    \item Solve
    \[
        K_k\Delta\mu_k=-R_k
    \]
    and update
    \[
        \mu^{(k+1)}=\mu^{(k)}+\Delta\mu_k.
    \]

    \item Run the integrator again with $\mu^{(k+1)}$, evaluate $R_{k+1}$,
    and compute
    \[
        s_k=\mu^{(k+1)}-\mu^{(k)},
        \qquad
        y_k=R_{k+1}-R_k.
    \]

    \item Update
    \[
        K_{k+1}
        =K_k+
        \frac{(y_k-K_ks_k)s_k^{\mathsf T}}
             {s_k^{\mathsf T}s_k}
    \]
    and continue the iteration.
\end{enumerate}

After convergence, the physical state is reconstructed as
\[
    z_{n+1}
    =\mathcal P_n\bigl(Y_{n+1}(\mu_n),\mu_n\bigr).
\]

During all internal iterations, the state $z_n$ remains fixed. Each new value
of $\mu^{(k)}$ requires rerunning the integrator because both $Y_n(\mu)$ and
$Y_{n+1}(\mu)$ depend on the multiplier. It is useful to distinguish the
indices
\[
    \boxed{
    n=\text{time step},
    \qquad
    k=\text{internal Broyden iteration}.}
\]

\section{Specialization to Different Integrators}

\subsection{ABBA scheme with two copies}

In the usual ABBA formulation,
\[
    \mathcal E_n(z_n,\mu)
    =\begin{pmatrix}z_n+\mu\\z_n-\mu\end{pmatrix},
    \qquad
    \Psi_{h,n}=\Phi_{\mathrm{ABBA}},
\]
and, if
\[
    \begin{pmatrix}u^{(2)}(\mu)\\v^{(2)}(\mu)\end{pmatrix}
    =\Phi_{\mathrm{ABBA}}
     \begin{pmatrix}z_n+\mu\\z_n-\mu\end{pmatrix},
\]
the compatibility condition is
\[
    \mathcal C_n
    \left(
        \begin{pmatrix}u^{(2)}\\v^{(2)}\end{pmatrix},
        \mu
    \right)
    =u^{(2)}-v^{(2)}+2\mu.
\]
Therefore,
\[
    R_n(\mu)=u^{(2)}(\mu)-v^{(2)}(\mu)+2\mu.
\]
When $h=0$, $\Phi_{\mathrm{ABBA}}$ is the identity and
\[
    R_n(\mu)=4\mu,
    \qquad
    K_0=4I_q
\]
is a natural initial approximation for small $h$. After convergence,
\[
    z_{n+1}=u^{(2)}(\mu_n)+\mu_n
           =v^{(2)}(\mu_n)-\mu_n.
\]

\subsection{BM4 scheme}

To use Broyden's method with BM4, it is enough to replace the base integrator
by
\[
    \Psi_{h,n}=\Phi_{\mathrm{BM4}}.
\]
If BM4 uses the same two-copy lifting and the same compatibility condition,
the maps $\mathcal E_n$, $\mathcal C_n$, and $\mathcal P_n$ remain unchanged.
If the BM4 architecture uses different auxiliary variables or a different
final condition, only these maps need to be redefined. The iterations
\eqref{eq:linear-correction}--\eqref{eq:good-broyden} remain exactly the same.

\subsection{Any other method}

For a new integrator, it is only necessary to specify:
\begin{equation}
\boxed{
\begin{aligned}
    Y_n(\mu)      &=\mathcal E_n(z_n,\mu),\\
    Y_{n+1}(\mu)  &=\Psi_{h,n}\bigl(Y_n(\mu)\bigr),\\
    R_n(\mu)      &=\mathcal C_n\bigl(Y_{n+1}(\mu),\mu\bigr),\\
    z_{n+1}       &=\mathcal P_n\bigl(Y_{n+1}(\mu_n),\mu_n\bigr).
\end{aligned}}
\end{equation}
Once the function $R_n$ has been defined, Broyden's method acts as a nonlinear
solver independent of the integrator used to evaluate it.

\section{Conclusion}

The fundamental object for Broyden's method is not the integrator, but the
residual $R_n(\mu)$. ABBA, BM4, or any other method enters only through the
evaluation of this residual. Consequently, a modular implementation should
separate
\[
    \boxed{
    \text{evaluation of }R_n(\mu)
    \quad\text{and}\quad
    \text{solution of }R_n(\mu)=0\text{ using Broyden's method}.}
\]
This separation makes it possible to reuse the same Broyden solver with
different integrators without modifying its update formulas.

\end{document}
